\section{Local Weil--Deligne parameters and their Galois realizations}
\label{sec:local-global}

The lift construction specifies a finite type on the automorphic side. Later we will
need to recognize that same local datum in the Galois representation of an ordinary
constituent. We match the local cuspidal
and principal series data to Weil--Deligne parameters and
specify the global attachment and its hypotheses. The
ordinary-level arguments below use these compatibilities
at both the split ramified places and the coefficient prime. Proposition~\ref{prop:prescribed-local-cuspidal-types} gives the local parameter on the specified cuspidal types. Definition~\ref{hyp:local-global-ranges} specifies the places and twists at which a global representation must recover that parameter.

Let \(v\) be a finite place of a CM field \(E\).
The local assignment constructed here associates a Weil--Deligne parameter to each irreducible smooth representation of \(\mathrm{GL}_n(E_v)\). In rank one it is given by Artin reciprocity.
\begin{proposition}[Local parameters on the prescribed cuspidal types]\label{prop:prescribed-local-cuspidal-types}
At a finite place \(v\) the local assignment sends a smooth irreducible representation of \(\mathrm{GL}_n(E_v)\) to an \(n\)-dimensional Weil--Deligne parameter. On characters of \(E_v^\times\), use Artin reciprocity.
At \(E_v=\Q_p\), use induction from the degree-\(n\) unramified extension on each depth-zero cuspidal, with determinant fixed by its central character. For the regular residue character \(\theta\), the inertia representation is 
\[
 \bigoplus_{j=0}^{n-1}\theta^{p^j}.
\]
 Extend by the Artin twist to its smooth
twists, independently of the presentation. On the remaining
cuspidals use irreducible induced parameters with nonscalar
wild inertia, constructed using the first principal unit group. The full assignment extends by
the Langlands segment rule.
\end{proposition}

\begin{proof}
\proofstep{The rank-one and depth-zero normalizations.}

At a rational split completion \(E_v=\Q_p\) the other specified
cuspidals are twists of depth-zero cuspidals. The latter
are compact inductions from the finite cuspidal types of
\(\mathrm{GL}_n(\mathbf F_p)\), extended to the group generated
by the maximal compact and the center. Their finite types are
indexed by Frobenius orbits of regular characters \(\theta\) of
\(\mathbf F_{p^n}^{\times}\). The parameter is induced from the
Weil group of the degree-\(n\) unramified extension. On inertia
it is the direct sum of the characters \(\theta^{p^j}\) for
\(0\le j<n\).

Write \(F\) for the chosen Frobenius lift and \(\gamma\)
for the value of the inducing character at \(F^n\).
In a cyclic basis \(F e_j=e_{j+1}\) for \(j<n-1\)
and \(F e_{n-1}=\gamma e_0\). Thus
\[
        \det(F)=(-1)^{n-1}\gamma .
\]
For central value \(b\) at the uniformizer in the reciprocity
convention, choose \(\gamma=(-1)^{n-1}b\). This is the
depth-zero normalization of the local Langlands correspondence.
\hfill\mbox{\cite{adrian-roe}}

\proofstep{Independence of twisting presentations and segments.}

Twisting by a smooth character of \(E_v^{\times}\) twists the
parameter by its Artin image. Suppose two depth-zero cuspidals differ by a twist \(\alpha\circ\det\). Their usual Weil parameters are tame. For wild inertia \(P_v\), compatibility of the usual correspondence with twists gives
\[
 I_n=\operatorname{rec}(\pi\otimes(\alpha\circ\det))(g)
   =\operatorname{rec}_1(\alpha)(g)I_n\qquad(g\in P_v).
\]
Thus the relative character is trivial on \(1+p\mathbb Z_p\) by local reciprocity. \cite{henniart-local,adrian-roe} The corresponding
finite cuspidal types differ by the determinant character
modulo \(p\).

On the elliptic torus of the finite group, multiplication
by \(u\in\mathbf F_{p^n}^\times\) has determinant equal to
its field norm. Consequently the finite twist on the torus is
\[
       (\alpha\circ\det)(u)
          =\alpha\bigl(N_{\mathbf F_{p^n}/\mathbf F_p}(u)\bigr).
\]
Functoriality of local reciprocity gives the same pullback
on the degree-\(n\) Weil group. Thus the inertial orbits
agree in the two presentations. The central determinant gives
the same last Frobenius coefficient in the cyclic basis, which
identifies the two induced parameters.
For a general irreducible based on these cuspidals, use the unique multisegment of its Langlands quotient. For segment lengths \(s_i\) and cuspidal parameters \(W_i\) the contribution is
\[
 \bigoplus_i\bigl(W_i\otimes\operatorname{Sp}(s_i)\bigr)\otimes|\cdot|^{u_i},
\]
where \(u_i\) is the central exponent of the segment in the chosen normalization. The multisegment classification ensures independence of their order. \cite[\S9]{zelevinsky} This preserves the
dimension and determines the monodromy lengths.

\proofstep{Auxiliary parameters on the remaining cuspidals.}

For the remaining cuspidals of dimension
\(n\ge2\), choose auxiliary parameters as follows. Over the unramified extension
of degree \(n\), the first principal unit quotient gives
the additive group of its residue field. Write \(k_v\) for the residue field of \(E_v\),
and \(k_n/k_v\) for the degree-\(n\) extension. Choose \(b\in k_n\)
of degree \(n\) and use the additive character
\[
    x\longmapsto
       \exp(2\pi i\,\operatorname{Tr}_{k_n/\mathbf F_p}(bx)/p)
\]
on the first principal unit quotient. Identify \(\mathbf F_p\)
with residues in the exponent.

To extend this character to the multiplicative group, write \(F_n/E_v\) for the unramified extension, \(\varpi\) for a uniformizer in \(E_v\), and \(U_n^i=1+\varpi^i\mathcal O_{F_n}\). Reduction and Teichm\"uller lifts give
\[
 F_n^\times/U_n^2
   =\varpi^{\mathbb Z}\times [k_n^\times]\times (U_n^1/U_n^2),
 \qquad 1+\varpi x\longmapsto\bar x\quad\text{on }U_n^1/U_n^2.
\]

Extend trivially across the first two factors. Local reciprocity then gives the inducing character of the Weil group. Write \(q_v=|k_v|\) and \(\lambda_c(x)=\exp(2\pi i\operatorname{Tr}(cx)/p)\).
Nondegeneracy of the trace pairing and the degree of \(b\) give
\[
 \lambda_{b^{q_v^i}}=\lambda_{b^{q_v^j}}
 \quad\Longleftrightarrow\quad b^{q_v^i}=b^{q_v^j}
 \quad\Longleftrightarrow\quad i\equiv j\pmod n .
\]
The inducing character thus gives \(n\) distinct inertia lines. Every invariant subspace for inertia
is a sum of these lines, and Frobenius permutes them transitively,
so the induced parameter is irreducible. By local reciprocity their
restriction to wild inertia has a nonscalar element, since
the additive quotient is a \(p\)-group and the tame quotient
has prime-to-\(p\) finite quotients. 

For a cuspidal \(\tau\) of rank \(n\), write \(W_0\) for this induction. In the convention \(|\varpi|_v=q_v^{-1}\) there is a unique \(t\in\mathbb R\) such that
\[
 q_v^{-nt}|\det W_0(\operatorname{Art}_v(\varpi))|
       =|\omega_\tau(\varpi)|.
\]
Use \(W_0\otimes |\cdot|^t\). Twisting \(\tau\) by \(|\det|^u\) changes \(t\) to \(t+u\), leaving its wild inertia action unchanged.

It remains to check that this auxiliary choice cannot obscure the parameters on the specified class. On a character, or on a twist of a depth-zero cuspidal, wild inertia acts by scalars on the irreducible Weil factor. For the auxiliary induction it acts by at least two distinct characters. The uniqueness of the decomposition of a Frobenius-semisimple parameter into indecomposable summands therefore separates the two cases:
\[
 W\otimes\operatorname{Sp}(s)\simeq W'\otimes\operatorname{Sp}(s')
       \ \Longrightarrow\ s=s',\quad W\simeq W'.
\]
Indeed \(s\) is the common Jordan block length and the action on the kernel of monodromy recovers \(W\) up to the prescribed norm twist. When every constituent belongs to the specified class, the construction coincides with the ordinary local correspondence by the depth-zero formula and compatibility with twists and segments. Injectivity there follows from the local Langlands bijection. This proves the separation property used below. \cite{henniart-local,adrian-roe,zelevinsky}

\end{proof}

\subsection{Agreement for every admissible map}

It remains to see why the later local calculations do not depend on a choice among
parameter maps satisfying the comparison properties. We call a local parameter map \(R\) admissible when it preserves
dimension, satisfies the rank-one Artin formula, assigns
irreducible parameters to cuspidals, and extends by the segment
rule above. It respects real norm twists and gives the absolute
value of the central character as the absolute determinant
under reciprocity.

On representations whose cuspidal support consists of
characters and, at rational completions, twists of depth-zero
cuspidals, it also respects smooth twists and contragredients and gives
\[
        \det R(\sigma)\circ\operatorname{Art}_v
            =\omega_\sigma ,
\]
where \(\omega_\sigma\) is the central character in the fixed
reciprocity convention. It gives tame parameters on this class
when there are first congruence invariants. If one representation belongs to this class,
equality of its parameter with that of any irreducible
implies equality of the representations. At a rational completion, if the restriction to
\(\mathrm{GL}_n(\mathbb Z_p)\) contains the inflation of the
regular finite cuspidal type, the parameter satisfies the
zero-monodromy and inertia-polynomial conditions of Proposition~\ref{prop:prescribed-local-cuspidal-types}.

For the inertia formula use reciprocity on the unramified extension that realizes the residue character. The determinant identity, the separation of cuspidal parameters by wild inertia and the segment construction in Proposition~\ref{prop:prescribed-local-cuspidal-types} exhibit such a map. On the specified class contragredients and smooth twists agree with those of the ordinary correspondence. \cite{henniart-local,adrian-roe}

\begin{proposition}[Agreement in the two local comparison ranges]
Let \(R\) be an admissible map. If \(R(\sigma)\) has zero monodromy
and is a sum of lines, the rank-one Artin formula and the segment
rule determine it. At a rational split completion, if \(R(\sigma)\)
is tame with zero monodromy, the cuspidal support of \(\sigma\)
consists of characters and depth-zero cuspidals. In each of these two ranges, writing \(\operatorname{rec}_{\rm norm}\) for the usual normalized correspondence, 
\[
 R(\sigma)\simeq\operatorname{rec}_{\rm norm}(\sigma).
\]

\end{proposition}

\begin{proof}
\proofstep{Singleton segments and toral parameters.}

The global comparison is required for every admissible map.
Let \(R\) be one such map. The special parameter of a segment with cuspidal
parameter \(W\) and length \(s\) has \(\dim W\) monodromy
blocks of size \(s\). Thus in the toral comparison
\[
 N_{R(\sigma)}=0\ \Longrightarrow\ s=1\text{ on each segment},
 \qquad R(\sigma)\text{ a sum of lines}\ \Longrightarrow\ \dim W=1.
\]
The dimension law and rank-one Artin formula then determine the parameter.

\proofstep{Count tame cuspidal parameters of fixed determinant.}

At a rational split completion, suppose instead that
\(R(\sigma)\) is tame with zero monodromy. We first show
that an irreducible rank-\(r\) cuspidal \(\tau\) with tame
\(R(\tau)\) has depth zero. Fix a lift of Frobenius and
write \(b\ne0\) for its determinant on \(R(\tau)\).
In a tame irreducible Weil representation the inertia
characters form one Frobenius orbit. Each occurs once:
on one inertia eigenspace choose an eigenvector for the
Frobenius power that stabilizes it, and take the span of its
Frobenius translates. Irreducibility makes this span the
entire space. Thus the orbit has length \(r\) and is
labelled by a regular character of
\(\mathbf F_{p^r}^{\times}\), up to Frobenius.

In a cyclic basis the determinant calculation above gives
\[
      \det(XI-F)=X^r-\gamma,\qquad \gamma=(-1)^{r-1}b .
\] For each regular orbit, the
diagonal inertia action together with this cyclic Frobenius
matrix defines an irreducible tame parameter. Two such
parameters are isomorphic exactly when the inertia orbits
coincide, since rotating a cyclic basis keeps the same last
coefficient. This gives the claimed bijection at fixed \(b\).

\proofstep{Use admissible injectivity on the specified class.}

For each label there is a distinct depth-zero cuspidal
type, and compact induction with the indicated central
value gives a depth-zero cuspidal with central character
of value \(b\) at the image of the chosen Frobenius.
Let \(\mathcal D_b\) be the resulting depth-zero cuspidals
with determinant value \(b\), and \(\mathcal W_b\) the tame
irreducible parameters with that value. Their two constructions
use the same finite set of regular Frobenius orbits. The depth,
dimension and central character laws give
\[
          R:\mathcal D_b\lhook\joinrel\longrightarrow\mathcal W_b,
          \qquad |\mathcal D_b|=|\mathcal W_b|.
\]
The arrow is injective by the admissible injectivity law.
It is therefore onto.

Since \(R(\tau)\in\mathcal W_b\),
the same law, now with the other representation in
\(\mathcal D_b\), gives \(\tau\in\mathcal D_b\). Applying
this to each singleton segment puts the entire support of
\(\sigma\) in the specified class. On each of these cuspidals
the finite type law gives the inertia orbit and the central
character law fixes the Frobenius determinant. Hence the
parameter of \(\sigma\) agrees with the usual normalized
correspondence. This argument applies in particular when the
tame inertia polynomial is squarefree, as required by the
global row. The spherical condition above \(\ell\) is a
condition on the smooth local representation itself.
\hfill\mbox{\cite{adrian-roe}}\label{cor:cohomology-local-global-1}

\end{proof}

\begin{construction}
\begin{tabular}{@{}ll@{}}\toprule
Restriction on \(R(\sigma)\) & Consequence for cuspidal support\\\midrule
\(N=0\) & every segment has length one\\
\(N=0\), sum of lines & every cuspidal has rank one\\
\(N=0\), tame, at a rational split place &
characters and depth-zero cuspidals\\\bottomrule
\end{tabular}
\caption{The two local comparison ranges. The segment rule detects monodromy block sizes. For a tame irreducible factor the Frobenius orbit of its inertia character identifies its depth-zero type.}\label{fig:two-local-comparison-ranges}
\end{construction}

\begin{paramdefinition}[Local comparisons required of a single representation]\label{hyp:local-global-ranges}
Let \((\pi,\chi)\) be a regular algebraic, polarized cuspidal
pair of degree \(n\) over \(E\), and fix a prime \(\ell\) and
an identification \(\iota:\overline{\Q}_\ell\simeq\C\).
For a continuous semisimple representation
\(r\) of dimension \(n\), define the following comparison with the normalized local parameters \(W_v\) of \(\pi_v\). Away from \(\ell\) require equality of dimensions, and in
the zero-monodromy toral case compare the full parameter after
Frobenius semisimplification. The second comparison concerns
squarefree tame inertia:
\[
  E_v=\Q_p,\qquad p>16n^2\ \text{if }n>1,
  \qquad
  \det(1-X\operatorname{WD}(r|_{G_{E_v}})(i))
     =\det(1-XW_v(i))\quad(i\in I_v),
\]
with zero monodromy. Above \(\ell\), a spherical local
representation requires crystallinity and the prescribed
Hodge--Tate weights. The same toral and tame cases above \(\ell\)
require potential semistability, the Hodge--Tate weights
and the indicated Weil--Deligne comparisons. The comparison holds for \(r\) if these conditions hold at every place in the indicated ranges.
\end{paramdefinition}

The agreement just proved transfers the comparisons
for the usual correspondence to every admissible map at
the places where these comparisons are requested.

\begin{proposition}[Simultaneous Galois attachment]\label{prop:simultaneous-local-global-attachment}
For any pair and coefficient prime in
Definition~\ref{hyp:local-global-ranges}, there exists one continuous semisimple representation 
\[
 r:G_E\longrightarrow\mathrm{GL}_n(\overline{\mathbb Q}_\ell)
\]
 with the prescribed Hodge--Tate and crystallinity properties. For every admissible local map, this same representation satisfies the stated Weil--Deligne comparison at every place in its range.
\end{proposition}

\begin{proof}
\proofstep{The usual correspondence and the algebraic weights.}

Barnet-Lamb, Gee, Geraghty and Taylor give the representation
with the prescribed Hodge--Tate weights and compare its
Frobenius-semisimplified parameter away from
\(\ell\) with
\[
             \operatorname{rec}(\pi_v\otimes
                              |\det|_v^{(1-n)/2}).
\]
For a highest algebraic weight \((a_{\tau,1},\ldots,a_{\tau,n})\)
the entries of the infinitesimal character of the dual
algebraic representation, and the resulting Hodge--Tate
jumps, are
\[
 z_{\tau,j}=-a_{\tau,j}-\frac{n+1-2j}{2},
 \qquad
 \frac{n-1}{2}-z_{\tau,j}=a_{\tau,j}+n-j .
\] Purity at the two conjugate embeddings
makes their multisets complementary with respect to the
purity weight. Their Theorem~2.1.1 gives de Rhamness and
crystallinity at spherical places over \(\ell\).
\hfill\mbox{\cite{blggt}}

\proofstep{Potential semistability and the monodromy partition.}

At other places over \(\ell\) the de Rham representation is
potentially semistable by the \(p\)-adic monodromy theorem.
Theorem~1.1 of the local-global compatibility theorem gives
the semisimple Weil action and bounds the monodromy
partition by that of the normalized local parameter in
each inertial class. In a nonempty inertial class, let \(\lambda_1\) be the largest
actual Jordan block and \(\mu_1\) the largest block of the bounding
parameter. The first dominance inequality gives
\[
    \lambda_1\le\mu_1=1
      \qquad\text{if the bounding monodromy vanishes}.
\]
Thus every block has size one, and monodromy is zero. Therefore in the toral and
tame situations just described the comparison gives the
required zero monodromy even at a regular weight with no
additional regularity condition. The rank-one case also
follows from the Galois character associated to the
algebraic Hecke character by global reciprocity.
\hfill\mbox{\cite{blggt-lg}}

\end{proof}

This identifies the local Galois data that enter the cohomology calculation.
Uniformity as the level, weights and rank vary will come from the length estimates
there.
